\section{Conclusion}
\label{sec:conclusion}

This paper proposes frequent losers---firms that never win yet
participate in abnormally many public tenders---as a screening marker for
cover bidding in procurement auctions. Using 4.5 million tender-items
from Brazil's BEC platform (2009--2019), we find strong support for
FL as a screening marker (H1): FL-present tenders exhibit
3.6--6.6\% higher negotiated prices, FL firms co-participate with
CADE-convicted cartelists at 3.5 times the expected rate
($p < 0.001$), and the FL screen achieves AUC~$= 0.94$---dominating
Imhof-style bid-level screens (AUC~$= 0.79$) while requiring only
participation data.

The evidence is also consistent with cover bidding as the causal
mechanism (H2): the FL--price effect concentrates in competitive
markets (Proposition~\ref{prop:market_selection}), Bajari--Ye tests
reject exchangeability and conditional independence, and the
structural estimation selects Regime~2 (coordinated cover bidding,
$\hat{\sigma}_c / \hat{\sigma}_g = 0.72$). Counterfactual analysis
shows that removing the convite minimum-bidder rule would save
R\$74--211 million, and the welfare-maximizing screening threshold
(1.6$\times$IQR at R\$100K per investigation) is close to the
paper's baseline (1.5$\times$).

\paragraph{Caveats.}
The staggered DiD is underpowered; the FL definition necessarily
involves misclassification; and strategic adaptation (rotating cover
bidders, allowing occasional wins) is a concern for operational
deployment.

\paragraph{Policy implications.}
The FL screen offers three advantages: (i)~requires only participation
and outcome data; (ii)~produces firm-level flags enabling targeted
investigation; (iii)~complements bid-level screens in a two-stage
workflow. The estimated welfare loss ranges from R\$400 million
(cross-fit) to R\$2.4 billion (IV) across 2009--2019. FL flags are
a starting point for investigation, not standalone enforcement evidence.

\paragraph{Future research.}
Four directions merit investigation: external validation beyond BEC;
deeper network analysis (shared representatives, financial links);
dynamic models of FL-screen--leniency interaction; and field
experiments in which procurement agencies adopt FL screening.
